% ============================================================================
% CANONICAL BASELINE PROMPT SPECIFICATION (π_formal)
% ============================================================================
% Status: FINAL - Do not modify without theoretical justification
% Date: 2025-12-28
% Purpose: Controlled reference condition for α measurement
% ============================================================================

\documentclass[11pt,a4paper]{article}
\usepackage[utf8]{inputenc}
\usepackage[T1]{fontenc}
\usepackage{amsmath,amssymb,amsthm}
\usepackage{booktabs}
\usepackage{tcolorbox}
\usepackage[margin=1in]{geometry}

\newtcolorbox{warningbox}{
  colback=red!5!white,
  colframe=red!75!black,
  fonttitle=\bfseries,
  title=DO NOT MODIFY,
  sharp corners
}

\newtcolorbox{specbox}{
  colback=blue!5!white,
  colframe=blue!50!black,
  sharp corners
}

\title{\textbf{Canonical Baseline Prompt Specification}\\
\Large $\pi_{\text{formal}}$ --- Controlled Reference Condition}
\author{FehrAdvice \& Partners AG\\BEATRIX Research Group}
\date{December 2025 --- FINAL}

\begin{document}
\maketitle

% ============================================================================
\section{Theoretical Status}
% ============================================================================

\begin{warningbox}
This prompt is the \textbf{controlled reference condition} from which all intervention effects are measured. Any modification invalidates comparability.
\end{warningbox}

\subsection{What This Prompt IS}

\begin{itemize}
    \item[$\checkmark$] Maximally unambiguous
    \item[$\checkmark$] Minimizes interpretation variance
    \item[$\checkmark$] The null point for effect measurement
    \item[$\checkmark$] Theoretically grounded specification
\end{itemize}

\subsection{What This Prompt Is NOT}

\begin{itemize}
    \item[$\times$] NOT an activation prompt
    \item[$\times$] NOT optimized for high $\alpha$
    \item[$\times$] NOT a conversational interaction
\end{itemize}

\subsection{Formal Definition}

\begin{equation}
\pi_{\text{formal}} := \arg\min_{\pi} \text{Amb}(\pi)
\end{equation}

where $\text{Amb}(\pi)$ denotes the interpretive ambiguity of prompt $\pi$.


% ============================================================================
\section{The Canonical Prompt}
% ============================================================================

\begin{specbox}
\textbf{Aufgabe}

Gegeben ist eine endliche Menge von Effekten:
\[
E = \{e_1, e_2, \dots, e_n\}
\]

\textbf{Ziel}

Ordne jedem Effekt $e_i$ genau eine Kategorie aus der Menge
\[
C = \{\text{REAL},\ \text{UNSICHER},\ \text{UNBEKANNT}\}
\]
zu.

\textbf{Domänenbedingung}

Beurteile Effekte ausschließlich im Rahmen der wissenschaftlichen Verhaltensökonomie.

\textbf{Epistemische Regel}
\begin{itemize}
    \item Wähle \textbf{REAL}, nur wenn du dich an konkrete, etablierte Fachliteratur (z.B. bekannte Effekte, Autoren oder Studien) erinnern kannst.
    \item Wähle \textbf{UNSICHER}, wenn der Effekt plausibel erscheint, du aber keine eindeutige Evidenz erinnern kannst.
    \item Wähle \textbf{UNBEKANNT}, wenn dir keine wissenschaftliche Evidenz erinnerbar ist.
\end{itemize}

\textbf{Prozedur}

Für jedes $e_i \in E$:
\begin{enumerate}
    \item Prüfe, ob dir wissenschaftliche Evidenz aus der Verhaltensökonomie bekannt ist.
    \item Ordne genau eine Kategorie aus $C$ zu.
    \item Triff keine zusätzlichen Annahmen.
\end{enumerate}

\textbf{Ausgabeformat}

Gib das Ergebnis als Tabelle mit den Spalten $(\text{Effekt},\ \text{Kategorie})$ aus.

\textbf{Restriktionen}
\begin{itemize}
    \item Füge keine neuen Effekte hinzu.
    \item Verwende keine externen Quellen.
    \item Erfinde keine Evidenz.
    \item Wenn keine Erinnerung an Fachliteratur vorliegt, wähle UNBEKANNT.
\end{itemize}
\end{specbox}


% ============================================================================
\section{Structural Analysis}
% ============================================================================

\subsection{Why This Prompt Is Formally Valid}

\begin{table}[h]
\centering
\caption{Formal requirements checklist}
\begin{tabular}{lcc}
\toprule
\textbf{Criterion} & \textbf{Fulfilled} & \textbf{Location in Prompt} \\
\midrule
Explicit goal & $\checkmark$ & ``Ordne ... genau eine Kategorie zu'' \\
Explicit decision space & $\checkmark$ & $C = \{\dots\}$ \\
Domain binding & $\checkmark$ & ``ausschließlich Verhaltensökonomie'' \\
Epistemic condition & $\checkmark$ & REAL only with recallable literature \\
Procedure & $\checkmark$ & 3 explicit steps \\
Quantifiers & $\checkmark$ & ``für jedes $e_i \in E$'' \\
No knowledge injection & $\checkmark$ & No examples, no effects provided \\
Minimal ambiguity & $\checkmark$ & No implicit defaults \\
\bottomrule
\end{tabular}
\end{table}

\subsection{Key Insight}

\begin{quote}
\textbf{This is a specification, not a conversation.}
\end{quote}


% ============================================================================
\section{Invariants (DO NOT VIOLATE)}
% ============================================================================

\begin{warningbox}
The following modifications \textbf{break theoretical validity}:

\begin{enumerate}
    \item[$\times$] Adding examples (injects knowledge, changes $M_{\text{latent}}$)
    \item[$\times$] Phrases like ``typically'', ``often'', ``it is known'' (introduces priors)
    \item[$\times$] Motivational phrases like ``be careful'' (activates $\alpha$, no longer baseline)
    \item[$\times$] Implicit evaluation standards (changes decision boundary)
    \item[$\times$] Changing the category set $C$ (incomparable results)
    \item[$\times$] Removing the domain constraint (unbounded task)
\end{enumerate}
\end{warningbox}


% ============================================================================
\section{Experimental Usage}
% ============================================================================

\subsection{Role in Experimental Design}

\begin{table}[h]
\centering
\begin{tabular}{ll}
\toprule
\textbf{Component} & \textbf{Specification} \\
\midrule
This prompt & Baseline condition \\
All other prompts & Interventions (relative to baseline) \\
Measured quantities & $\mathbb{E}[\alpha_{\text{BE}}]$, $\text{Var}(\alpha_{\text{BE}})$ \\
\bottomrule
\end{tabular}
\end{table}

\subsection{Interpretation of Results}

For any intervention prompt $\pi'$:

\begin{equation}
\Delta\alpha = \alpha(\pi') - \alpha(\pi_{\text{formal}})
\end{equation}

\begin{itemize}
    \item $\Delta\alpha > 0$: Intervention increases activation competence
    \item $\Delta\alpha = 0$: No effect
    \item $\Delta\alpha < 0$: Intervention decreases activation competence (unlikely but testable)
\end{itemize}


% ============================================================================
\section{Copy-Paste Ready Version}
% ============================================================================

\begin{verbatim}
**Aufgabe**

Gegeben ist eine endliche Menge von Effekten:
E = {e_1, e_2, ..., e_n}

**Ziel**

Ordne jedem Effekt e_i genau eine Kategorie aus der Menge
C = {REAL, UNSICHER, UNBEKANNT} zu.

**Domänenbedingung**

Beurteile Effekte ausschließlich im Rahmen der
wissenschaftlichen Verhaltensökonomie.

**Epistemische Regel**

- Wähle REAL, nur wenn du dich an konkrete, etablierte
  Fachliteratur (z.B. bekannte Effekte, Autoren oder Studien)
  erinnern kannst.
- Wähle UNSICHER, wenn der Effekt plausibel erscheint, du aber
  keine eindeutige Evidenz erinnern kannst.
- Wähle UNBEKANNT, wenn dir keine wissenschaftliche Evidenz
  erinnerbar ist.

**Prozedur**

Für jedes e_i ∈ E:
1. Prüfe, ob dir wissenschaftliche Evidenz aus der
   Verhaltensökonomie bekannt ist.
2. Ordne genau eine Kategorie aus C zu.
3. Triff keine zusätzlichen Annahmen.

**Ausgabeformat**

Gib das Ergebnis als Tabelle mit den Spalten
(Effekt, Kategorie) aus.

**Restriktionen**

- Füge keine neuen Effekte hinzu.
- Verwende keine externen Quellen.
- Erfinde keine Evidenz.
- Wenn keine Erinnerung an Fachliteratur vorliegt,
  wähle UNBEKANNT.
\end{verbatim}


% ============================================================================
\section{Summary}
% ============================================================================

\begin{quote}
\textbf{The mathematically formal prompt is the controlled reference condition that guarantees maximum unambiguity with minimal epistemic steering.}
\end{quote}

\begin{equation}
\boxed{\pi_{\text{formal}} = \text{Baseline} \implies \forall \pi': \Delta\alpha(\pi') = \alpha(\pi') - \alpha(\pi_{\text{formal}})}
\end{equation}


\end{document}
